% Appendix I: competing accounts and further related work. Input by sections/appendix.tex.
\section{Competing accounts and further related work}\label{app:accounts}\label{app:discussion}

\paragraph{What four accounts predict (as of the earlier entries).}
\Cref{tab:accounts} sets four accounts against the observations of the discovery sample. Each, ours included, is contradicted somewhere. Two rows matter most: the binding swap, where our account's preregistered extension failed (H7, H9) and the lookback account is partly compatible, post hoc (one edit, one model); and I6, where at 24B less of the remap's effect was kept than predicted once its readers were ablated.

\begin{table}[!ht]
  \centering
  \caption{\textbf{What four accounts predict, as of the entries of the discovery sample.} DT: duplicate-token attention alone; LB: lookback addresses and payloads \citep{prakash2025lookbacks,wu2026record}; FF: an edit's channel is set by the format it was fit in; LM: later-mention readers (ours). $\checkmark$: predicted and observed; $\times$: predicted otherwise (refuted); --: silent; (ph): compatible, post hoc; $\checkmark$/$\times$: the attention is predicted, its presence without a read is not; $\checkmark$/(ph): the attention is predicted, its presence without a read is explained post hoc (the read is decided downstream). Binding-swap row: \cref{app:binding}; last two rows: \cref{app:stage7}. The final round's results are added in the text.}
  \label{tab:accounts}
  {\small\begin{tabular}{@{}p{0.54\columnwidth}cccc@{}}
\toprule
Observation (evidence) & DT & LB & FF & LM \\
\midrule
Hop~1: a sparse set of heads, whose core is duplicate-token heads, attends from the re-mentions to the writing token, key-matched, at every scale, also at 1.5B and 3B where $\mathrm{ID}_K \approx 0$ (G1--G3, H1--H4) & $\checkmark$/$\times$ & -- & -- & $\checkmark$ \\
The same words placed before the writing token give no key read (E1) & $\checkmark$ & $\times$ & -- & $\checkmark$ \\
The answer reads the option words, not the writing token (H5, G8b) & -- & $\times$ & -- & $\checkmark$ \\
Without a later mention the identity is copied through the answer's edge to the writing token (G8a, E4) & -- & $\checkmark$ & -- & $\checkmark$ \\
Removing the readers raises the value copy (G7a, H3) & -- & (ph) & -- & $\checkmark$ \\
A remap refit without later mentions shows the same crossover (F1--F4) & -- & -- & $\times$ & $\checkmark$ \\
Prakash et al.'s binding swap is key-dominant in every evaluable format ($\kappa$ 0.618, $\psi_V$ 0.435 without a later mention); later mentions shift $\kappa$ by $+$0.247 (H7--H10) & -- & (ph) & -- & $\times$ \\
\bottomrule
\end{tabular}
}
\end{table}

\paragraph{What the final round adds.}
Four results of the final entry bear on the accounts. (i)~The first hop's output: moved into the row of an option absent from the story (and out of B's row), the reader heads' write moves the answer there and is read by key; added alone it moves the answer \JdDoseOneQwenSeven{} and \JdDoseOneMistralSeven{} as much as the natural clamp (\cref{app:flag}), so a duplicate-token match writes something the answer uses (\cref{app:flag}); but structured controls also move the answer, and removing one direction per layer leaves about half of the key read, so the write is sufficient, not specific, and not necessary as one direction. Duplicate-token attention alone (DT) predicts the attention, not this write; it does not decide between LB and LM. (ii)~The sign of the key read follows the task: asked which option was not mentioned, the key read and the value read both turn negative, and the same frozen reader heads carry both questions' reads; in IOI's in-sentence re-mentions the key read is negative while the value read stays positive, applied at the answer's key read of the listed names, not by the answer-side heads ranked in lists (\cref{app:sign}). A negative use of duplicate detection is known \citep{wang2022ioi,mcdougall2023copy}; we do not present the polarity as a discovery. (iii)~The channel-ratio law, which would have made an intervention's apparent channel a property of the model's own read in each format, failed, even for steering vectors close to the natural edit (\cref{app:claim3}); so the LM column's ``\checkmark'' for the refit row describes the remap of the release and its refit, not edits in general. (iv)~On counterfactual reading-comprehension passages a later options list opens the key read, but behaviour follows the key only with lettered options (\cref{app:natural}).

\paragraph{Relation to lookback, address and payload accounts.}
\citet{prakash2025lookbacks} find, with causal mediation and causal abstraction, lookback mechanisms built on ordering IDs (binding, answer and visibility lookbacks) that track beliefs; \citet{wu2026record} find that transplanting an operation's span redirects which state is read (a routing record), and that the writing model reads mainly the address. In our terms, a lookback whose payload is the state's identity would read identity by copy, through the value. The exchange of \cref{app:prakash} shows that their state-token patch, which swaps ordering information rather than identity, reaches the answer mainly through the state tokens' keys, with a substantial value share without a later mention ($\psi_K$ = \PkPsiKBindNone{}, $\psi_V$ = \PkPsiVBindNone{} under \fmtNone{}); because one patch moves both, the decomposition does not separate an address from a payload, and it is one edit at one model. What is copied through the value without a later mention is the state's identity: an identity edit at the same positions is value-carried under \fmtNone{} ($\psi_V$ = \PkPsiVIdZeroNone{}), as the released remap of \citet{anonymous2026fitted} is (E4), and in the final round so is a steering vector at block 28 (J-C6, \cref{app:binding-j}). Later mentions add a second route for identity, a lookup in which a re-mentioned word finds the writing token through its key, as duplicate-token heads do \citep{wang2022ioi}; at two 7B models it runs through a sparse set of heads enriched for duplicate-token heads, and the answer then reads the re-mentioned words (\cref{app:heads}). On their material this route is visible for an identity edit at block 0 ($\kappa_{\mathrm{ID}}$ rises from \PkKappaIdZeroNone{} under \fmtNone{} to \PkKappaIdZeroOpt{} under \fmtOpt{}) but closes with depth, and the final round's screen found no depth with both it and a working binding swap in two further models (\cref{app:binding-j}), so whether it runs alongside a lookback at matched depth is untested. Which of the positional, lexical and reflexive mechanisms of \citet{gurarieh2025mixing} our lookup corresponds to, if any, we do not test. In our \textsc{qnames2} format, which names both candidate states, the binding swap shifts the margin through keys alone ($\kappa$ = \PkKappaBindQnamestwo{}) but flips no answer (exploratory). As for the role-swap control: the observer is the first character named in both runs, so the control cannot detect an ordering ID of the observer at the writing token.

\paragraph{Further related work.}
On heads and circuits: attention weights alone do not determine a head's output \citep{kobayashi2020norm}, induction can copy tokens or whole concepts \citep{feucht2025dualroute}, and circuit components are reused across tasks \citep{merullo2024reuse}. Further work on options and retrieval: specific heads select among options \citep{lieberum2023chinchilla,wiegreffe2025answer}, select-and-copy heads do so by query--key alignment \citep{tulchinskii2024wise}, and a winner is chosen before it is bound to a symbol \citep{wong2026decide}; shallower heads that read persuasive keywords write a one-dimensional option-routing feature into the option tokens, which redirects the attention of decision heads that copy the option they attend to \citep{sun2026persuaded}, the closest analogue of our two hops; belief and reality share a value slot read out under a router at the query \citep{steele2026routing}; answering with option symbols requires binding the options to the symbols \citep{robinson2023mcsb}; questions address facts by their order of mention \citep{zhou2026ordinal}, compatible with (not tested by) the key-dominant binding swap; value-slot routers \citep{steele2026spaces} and cached-KV edits \citep{bojieli2026notes} concern storage. On the validity of interventions: steering vectors are not identifiable from input--output behaviour \citep{venkatesh2026steering}; diagnostics include control tasks \citep{hewitt2019control}, benchmarks \citep{huang2024ravel,arora2024causalgym} and input partitions \citep{li2026bucketing}; the format can change what a steering gain identifies \citep{gao2026encodings}; models' detection and localisation of injected concept vectors relies on gate and router heads, and better localisation is associated with better alignment of the key and value changes the injection induces in their QK and OV computations \citep{zou2026introspection}. Activation patching \citep{vig2020causal} depends on the metric and the corruption \citep{zhang2023patching,heimersheim2024patching}, and precomputed KV caches substitute exactly for a joint pass \citep{pustovit2026packs}.
